%% kapa: exam=2 round=36 year=2025 period=1 subject=appraisal_practice subject_ko=감정평가실무 seq=2 number=2 points=30
%% source: Q-Net 감정평가사 기출문제 (cst00309) artlSeq=5239325 file=raw/1008/5239325/2025년 제36회 감정평가사 2차 감정평가실무.pdf
%% conversion: kordoc 4.24.51 -> tools/gen2.py (auto); review=proofread-gpt-sol; flags=none
\begin{longquestion}{【문제 2】}
감정평가사 甲은 상품 판매업을 하는 개인기업 A의 법인전환에 따른 영업권 감정평가를 의뢰받았다. 관련 법규 및 이론을 참고하고, 제시된 자료를 활용하여 다음 물음에 답하시오. (30점) \par
\medskip\noindent 물음 1) 대상기업 A의 영업관련 기업가치를 평가하시오. (20점) \par
\medskip\noindent 물음 2) 대상기업 A의 영업권 가치를 평가하시오. (10점) \par
< 자료 1 > 대상기업 개요 \par
{\small\centering
\resizebox{\ifdim\width>\linewidth\linewidth\else\width\fi}{!}{%
\begin{tabular}{|c|c|c|c|}
\hline
\makecell{상호} & \makecell{A} & \makecell{대표자} & \makecell{乙} \\
\hline
\makecell{개업일} & \makecell{2017.01.01.} & \makecell{사업의 종류} & \makecell{도매업, 소매업} \\
\hline
\makecell{종목} & \makecell{전자상거래업, 소매업} & \makecell{주요상품} & \makecell{각종 액세서리 등} \\
\hline
\end{tabular}}\par}

< 자료 2 > 주요 가정 \par
1. 대상기업의 특성 및 시장상황 등을 고려하여 고속성장기는 5년(1기∼5기)이고, 6기부터는 안정성장기로 가정함. 안정성장기의 영구성장율은 0 \%를 적용하며, 5기와 동일한 현금흐름이 지속되는 것으로 가정함 (단, 고속성장기 1기는 2025.01.01. ∼ 2025.12.31.임) \par
2. 대상기업의 결산일은 매년 말일이며, 현금흐름은 편의상 기말에 발생하는 것으로 가정함 \par
3. 대상기업의 현금흐름 추정시 비영업용 자산에 의한 수익, 비용은 제외된 것으로 가정함 \par
< 자료 3 > 재무상태표 및 손익계산서(2024. 12. 31. 현재) \par
1. 재무상태표 \par
(단위: 원) \par
{\small
\begin{longtable}{|l|r|r|}
\hline
\makecell{과목} & \multicolumn{2}{|c|}{\makecell{금액}} \\
\hline
\endhead
\makecell{자산} &  &  \\
\hline
\makecell{Ⅰ. 유동자산} &  & \makecell{178,000,000} \\
\hline
\makecell{(1) 당좌자산} &  & \makecell{118,000,000} \\
\hline
\makecell{현금 및 현금성 자산} & \makecell{18,000,000} &  \\
\hline
\makecell{매출채권} & \makecell{80,000,000} &  \\
\hline
\makecell{선급비용} & \makecell{20,000,000} &  \\
\hline
\makecell{(2) 재고자산} &  & \makecell{60,000,000} \\
\hline
\makecell{상품} & \makecell{60,000,000} &  \\
\hline
\makecell{Ⅱ. 비유동자산} &  & \makecell{159,000,000} \\
\hline
\makecell{(1) 투자자산} &  & \makecell{8,000,000} \\
\hline
\makecell{매도가능증권} & \makecell{8,000,000} &  \\
\hline
\makecell{(2) 유형자산} &  & \makecell{151,000,000} \\
\hline
\makecell{토지} & \makecell{100,000,000} &  \\
\hline
\makecell{차량운반구} & \makecell{60,000,000} &  \\
\hline
\makecell{감가상각누계액} & \makecell{(10,000,000)} &  \\
\hline
\makecell{비품} & \makecell{5,000,000} &  \\
\hline
\makecell{감가상각누계액} & \makecell{(4,000,000)} &  \\
\hline
\makecell{(3) 무형자산} &  &  \\
\hline
\makecell{(4) 기타 비유동자산} &  &  \\
\hline
\makecell{자산총계} &  & \makecell{337,000,000} \\
\hline
\makecell{부채} &  &  \\
\hline
\makecell{Ⅰ. 유동부채} &  & \makecell{155,000,000} \\
\hline
\makecell{매입채무} & \makecell{70,000,000} &  \\
\hline
\makecell{미지급금} & \makecell{36,000,000} &  \\
\hline
\makecell{예수금} & \makecell{31,000,000} &  \\
\hline
\makecell{단기차입금} & \makecell{18,000,000} &  \\
\hline
\makecell{Ⅱ. 비유동부채} &  & \makecell{62,000,000} \\
\hline
\makecell{장기차입금} & \makecell{62,000,000} &  \\
\hline
\makecell{부채총계} &  & \makecell{217,000,000} \\
\hline
\makecell{자본} &  &  \\
\hline
\makecell{Ⅰ. 자본금} &  & \makecell{120,000,000} \\
\hline
\makecell{자본금} & \makecell{120,000,000} &  \\
\hline
\makecell{자본총계} &  & \makecell{120,000,000} \\
\hline
\makecell{부채 및 자본 총계} &  & \makecell{337,000,000} \\
\hline
\end{longtable}
}

2. 손익계산서 \par
(단위: 원) \par
{\small
\begin{longtable}{|l|r|r|}
\hline
\makecell{과목} & \multicolumn{2}{|c|}{\makecell{금액}} \\
\hline
\endhead
\makecell{Ⅰ. 매출액} &  & \makecell{1,000,000,000} \\
\hline
\makecell{상품매출} & \makecell{1,000,000,000} &  \\
\hline
\makecell{Ⅱ. 매출원가} &  & \makecell{600,000,000} \\
\hline
\makecell{기초상품재고액} & \makecell{50,000,000} &  \\
\hline
\makecell{당기상품매입액} & \makecell{610,000,000} &  \\
\hline
\makecell{기말상품재고액} & \makecell{60,000,000} &  \\
\hline
\makecell{Ⅲ. 매출총이익} &  & \makecell{400,000,000} \\
\hline
\makecell{Ⅳ. 판매비와 관리비} &  & \makecell{176,000,000} \\
\hline
\makecell{직원급여} & \makecell{75,000,000} &  \\
\hline
\makecell{상여금} & \makecell{12,000,000} &  \\
\hline
\makecell{퇴직급여} & \makecell{8,000,000} &  \\
\hline
\makecell{복리후생비} & \makecell{14,000,000} &  \\
\hline
\makecell{여비교통비} & \makecell{5,000,000} &  \\
\hline
\makecell{접대비} & \makecell{7,000,000} &  \\
\hline
\makecell{통신비} & \makecell{1,000,000} &  \\
\hline
\makecell{세금과공과금} & \makecell{11,000,000} &  \\
\hline
\makecell{감가상각비} & \makecell{10,000,000} &  \\
\hline
\makecell{보험료} & \makecell{4,000,000} &  \\
\hline
\makecell{차량유지비} & \makecell{8,000,000} &  \\
\hline
\makecell{운반비} & \makecell{3,000,000} &  \\
\hline
\makecell{소모품비} & \makecell{10,000,000} &  \\
\hline
\makecell{지급수수료} & \makecell{8,000,000} &  \\
\hline
\makecell{Ⅴ. 영업이익} &  & \makecell{224,000,000} \\
\hline
\makecell{Ⅵ. 영업외 수익} &  & \makecell{2,000,000} \\
\hline
\makecell{잡이익} & \makecell{2,000,000} &  \\
\hline
\makecell{Ⅶ. 영업외 비용} &  & \makecell{1,000,000} \\
\hline
\makecell{잡손실} & \makecell{1,000,000} &  \\
\hline
\makecell{Ⅷ. 소득세차감전 이익} &  & \makecell{225,000,000} \\
\hline
\end{longtable}
}

< 자료 4 > 매출액 및 매출원가 관련 자료 \par
1. 매출액은 고속성장기에서는 과거 3년간의 매년 상승률 추세가 지속될 것으로 예측되며, 안정성장기에는 상승률 0\%로 가정함. 상승률은 백분율로 소수점 첫째자리에서 반올림하고, 매출액은 십만원 단위에서 반올림함 \par
(단위: 원) \par
{\small\centering
\resizebox{\ifdim\width>\linewidth\linewidth\else\width\fi}{!}{%
\begin{tabular}{|c|c|c|c|}
\hline
\makecell{구분} & \makecell{2021년} & \makecell{2022년} & \makecell{2023년} \\
\hline
\makecell{매출액} & \makecell{915,141,600} & \makecell{942,595,900} & \makecell{970,873,800} \\
\hline
\end{tabular}}\par}

2. 매출원가는 2024년의 매출원가율과 동일한 매출원가율을 적용하여 추정함 \par
< 자료 5 > 판매비와 관리비, 임차료, 대표자 급여(자가 노력비) 관련 자료 \par
1. 판매비와 관리비는 2024년의 판매비와 관리비 비율과 동일한 비율을 적용하여 추정하되, 판매비와 관리비 비율에는 임차료가 포함되고, 대표자 급여(자가 노력비)는 포함되지 아니함 \par
2. 대상기업은 특수관계자의 건물을 임차해서 영업에 사용하고 있어 임차료를 지급하지 않고 있으며, 적정한 임대차 조건은 2024.12.31. 현재 매월 임차료 2,000,000원과 보증금 20,000,000원으로 조사되었음 \par
3. 乙은 실질적으로 영업활동에 기여하고 있으며, 관리자급 직원의 급여 수준은 2024.12.31. 현재 연간 70,000,000원이고, 고속성장기 동안 매년 1,000,000원씩 증가됨 \par
< 자료 6 > 자본적지출 관련 자료 \par
감가상각비와 동일한 금액이 자본적 지출로 재투자되는 것으로 가정함 \par
< 자료 7 > 세금 관련 자료 \par
1. 계산의 편의상 지방소득세를 포함하여 소득세율은 33\%, 법인세율은 22\%를 적용함 (단, 누진세율은 적용하지 아니함) \par
2. 세금 계산시 제시된 자료 외의 감면 등은 없는 것으로 가정함 \par
< 자료 8 > 추가운전자본 관련 자료 \par
1. 추가운전자본은 운전자본소요율을 이용하여 산정함 \par
\[
\text{운전자본소요율}=\frac{1}{\text{매출채권회전율}}+\frac{1}{\text{재고자산회전율}}-\frac{1}{\text{매입채무회전율}}
\]
추가운전자본 = 매출액 증가분 $\times$ 운전자본소요율 \par
2. 운전자본소요율 계산 자료 \par
{\small\centering
\resizebox{\ifdim\width>\linewidth\linewidth\else\width\fi}{!}{%
\begin{tabular}{|c|c|c|c|}
\hline
\makecell{구분} & \makecell{매출채권회전율} & \makecell{재고자산회전율} & \makecell{매입채무회전율} \\
\hline
\makecell{회} & \makecell{10} & \makecell{10} & \makecell{20} \\
\hline
\end{tabular}}\par}

< 자료 9 > 할인율 관련 자료 \par
1. 할인율은 가중평균자본비용(WACC)을 사용하며, 백분율로 소수점 둘째자리에서 반올림함 \par
2. 대상기업의 자본구조는 제시된 2024.12.31. 현재의 재무상태표를 기준으로 결정함 \par
3. 자기자본비용은 자본자산가격결정모형(CAPM: Capital Asset Pricing Model)에 의하여 결정하며, 무위험자산의 수익률(Rf)은 국고채 수익률 등을 고려하여 4.0\%, 시장의 기대수익률(E(Rm))은 10.0\%를 적용함 \par
4. β계수는 유사업종에 속한 기업들의 β계수 등을 고려하여 1.1로 결정함 \par
5. 대상기업의 규모 등을 고려한 위험프리미엄은 7.4\%임 \par
6. 타인자본비용 결정시 차입이자율은 시장의 대출금리 등을 종합적으로 고려하여 8.2\%를 적용함 \par
7. 고속성장기와 안정성장기의 가중평균자본비용(WACC)은 동일하다고 가정함 \par
< 자료 10 > 재무상태표 관련 자료 \par
1. 재무상태표상 비영업용 항목은 매도가능증권임 \par
2. 토지는 2017년에 취득한 영업 관련 자산으로, 현재가치(시가) 산정을 위해 감정평가한 결과는 160,000,000원임 \par
3. 그 외 자산, 부채는 장부가액과 시가와의 차이가 없음 \par
< 자료 11 > 기타 사항 \par
1. 기준시점: 2025.01.01. \par
2. 영업권은 대상기업의 영업관련 기업가치에서 영업투하자본을 차감하는 방법으로 평가할 것 \par
3. 영업관련 기업가치는 십만원 단위에서 반올림할 것 \par
\end{longquestion}
